\documentclass[11pt,a4paper]{article}
\usepackage[utf8]{inputenc}
\usepackage[T1]{fontenc}
\usepackage[french]{babel}
\usepackage{lmodern,geometry,microtype}
\usepackage{amsmath}
\usepackage{siunitx}
\usepackage[version=4]{mhchem}
\usepackage{xcolor,listings,hyperref}
\geometry{margin=2.4cm}
\sisetup{locale=FR,per-mode=symbol}
\DeclareSIUnit\molecule{molécule}
\DeclareSIUnit\ppm{ppm}
\lstset{language=Python,basicstyle=\ttfamily\small,backgroundcolor=\color{black!4},frame=single,breaklines=true,showstringspaces=false}
\hypersetup{colorlinks=true,linkcolor=blue!45!black,urlcolor=blue!55!black}

\title{Documentation de \texttt{code\_final\_emissivite\_avec\_sections\_hitran.py}}
\author{Projet climat 2026}
\date{\today}

\begin{document}
\maketitle
\tableofcontents

\section{Rôle du programme}

Ce module calcule l'épaisseur optique et l'émissivité spectrale d'une couche
atmosphérique contenant \ce{CO2}, \ce{H2O}, \ce{CH4} et \ce{O3}. Il combine
les sections efficaces HITRAN, un profil vertical de pression, température et
concentrations, puis la loi de Beer--Lambert.

\section{Installation et fichiers indispensables}

\begin{lstlisting}[language=bash]
python3.12 -m venv .venv
source .venv/bin/activate
python -m pip install -r requirements.txt
\end{lstlisting}

Le paquet externe directement importé est \texttt{numpy}. HAPI est requis
indirectement par \path{2_sections_efficaces_gaz_hitran.py}. Deux fichiers
locaux sont obligatoires :

\begin{itemize}
  \item \path{2_sections_efficaces_gaz_hitran.py}, placé dans le même dossier ;
  \item \path{../../temperatures_et_concentrations/code_H20_CH4_CO2_O3.py},
        placé dans l'arborescence actuelle du projet.
\end{itemize}

Le cache \path{donnees_hitran/} est nécessaire, ou doit pouvoir être rempli
depuis le serveur HITRAN.

\section{Imports et chargements dynamiques}

\begin{lstlisting}
from importlib import import_module
from pathlib import Path

import numpy as np
\end{lstlisting}

\texttt{importlib} et \texttt{pathlib} appartiennent à la bibliothèque
standard. Les deux dépendances locales sont ensuite chargées ainsi :

\begin{lstlisting}
module_sections_hitran = import_module(
    "2_sections_efficaces_gaz_hitran"
)
calculer_sections_gaz = (
    module_sections_hitran.calculer_sections_gaz
)

espace_code_concentrations = {}
exec(chemin.read_text(encoding="utf-8"),
     espace_code_concentrations)
get_gases_ppm = espace_code_concentrations["get_gases_ppm"]
\end{lstlisting}

Le premier chargement est nécessaire parce que le nom du module commence par
un chiffre. Le second exécute le fichier du profil atmosphérique ; ce fichier
doit donc être fiable. Le chemin est calculé depuis \texttt{\_\_file\_\_} et non
depuis le dossier courant.

\section{Grandeurs et unités attendues}

\begin{center}
\begin{tabular}{lll}
Grandeur & Argument & Unité \\
\hline
Pression & \texttt{pression} & \si{\pascal} \\
Température & \texttt{temperature} & \si{\kelvin} \\
Épaisseur de couche & \texttt{epaisseur\_couche} & \si{\metre} \\
Altitude & \texttt{altitude\_couche\_km} & \si{\kilo\metre} \\
Longueur d'onde & \texttt{longueur\_onde} & \si{\metre} \\
Concentration & valeurs du profil & \si{\ppm} \\
Section efficace & tableaux HITRAN & \si{\metre\squared\per\molecule}
\end{tabular}
\end{center}

\section{Fonctions élémentaires}

\subsection{Densité moléculaire de l'air}

\texttt{calculer\_densite\_moleculaire\_air(pression, temperature)} renvoie
\begin{equation}
  n_{\mathrm{air}}=\frac{p}{k_{\mathrm B}T},\qquad
  k_{\mathrm B}=\SI{1.380649e-23}{\joule\per\kelvin},
\end{equation}
en molécules par \si{\metre\cubed}.

\subsection{Conversion des concentrations}

\texttt{convertir\_concentrations\_en\_fractions\_molaires(concentrations)}
lit les clés \texttt{CO2/H2O/CH4/O3 de mélange (Proportion)}. Chaque valeur
en ppm est multipliée par $10^{-6}$. Le retour est un dictionnaire indexé par
\texttt{CO2}, \texttt{H2O}, \texttt{CH4} et \texttt{O3}.

\subsection{Densité d'un gaz}

\texttt{calculer\_densite\_moleculaire\_gaz(p, T, x)} renvoie
\begin{equation}
  n_g=x_g n_{\mathrm{air}}.
\end{equation}

\subsection{Épaisseur optique}

\texttt{calculer\_epaisseur\_optique\_gaz(...)} reçoit la pression, la
température, la fraction molaire, l'épaisseur de couche et la section efficace.
Elle renvoie la grandeur sans dimension
\begin{equation}
  \tau_g=\sigma_g n_g\Delta z.
\end{equation}

\subsection{Émissivité}

\texttt{calculer\_emissivite(epaisseur\_optique)} applique
\begin{equation}
  \varepsilon_\lambda=1-\exp(-\tau_\lambda).
\end{equation}
La fonction accepte un scalaire ou un tableau NumPy.

\section{Fonction principale : une couche}

\texttt{calculer\_emissivite\_une\_couche(pression, temperature,
epaisseur\_couche, altitude\_couche\_km, longueur\_onde)} suit ces étapes :

\begin{enumerate}
  \item appel de \texttt{get\_gases\_ppm} à l'altitude demandée ;
  \item conversion de la pression de Pa en atm pour HAPI ;
  \item calcul des quatre sections efficaces à la pression et à la
        température locales ;
  \item conversion des ppm en fractions molaires ;
  \item calcul des quatre épaisseurs optiques ;
  \item somme $\tau_{\mathrm{total}}=\sum_g\tau_g$ et calcul de l'émissivité.
\end{enumerate}

Le dictionnaire retourné contient \texttt{emissivite}, \texttt{tau\_total},
\texttt{tau\_CO2}, \texttt{tau\_H2O}, \texttt{tau\_CH4}, \texttt{tau\_O3},
\texttt{sections\_efficaces}, \texttt{concentrations} et
\texttt{fractions\_molaires}. Les valeurs sont scalaires ou vectorielles selon
la longueur d'onde fournie.

\section{Test autonome}

\begin{lstlisting}[language=bash]
python code_final_emissivite_avec_sections_hitran.py
\end{lstlisting}

Le bloc principal teste une couche de \SI{1}{\metre} au sol, à
\SI{101325}{\pascal}, \SI{288}{\kelvin} et \SI{10}{\micro\metre}. Il affiche
les sections, concentrations, fractions molaires, contributions optiques et
l'émissivité finale.

\section{Hypothèses et précautions}

Le calcul décrit une couche homogène et isotherme, sans diffusion, nuages ni
aérosols. L'égalité entre absorptivité et émissivité suppose l'équilibre
thermodynamique local. Une colonne non isotherme ne s'obtient pas en sommant
directement les émissivités : il faut ajouter l'émission de Planck de chaque
couche et son atténuation par les couches supérieures.

\end{document}
